\textbf{Phần 1. Trắc nghiệm lựa chọn}
\setcounterpageref{ex}{0}

\begin{ex} %[3H3V52]
Một hình chữ nhật có chiều dài $12\text{ cm}$ và chiều rộng $7\text{ cm}$. Diện tích của hình chữ nhật đó là
\choice
{$38\text{ cm}^2$}
{$84\text{ cm}$}
{$48\text{ cm}^2$}
{\True $84\text{ cm}^2$}
\loigiai{
Diện tích hình chữ nhật: $12 \times 7 = 84\text{ (cm}^2)$.
\textbf{Đáp án của câu hỏi là D}
}
\end{ex}

\begin{ex} %[3H3V52]
Một hình vuông có cạnh dài $9\text{ cm}$. Diện tích của hình vuông đó là
\choice
{$36\text{ cm}^2$}
{\True $81\text{ cm}^2$}
{$72\text{ cm}^2$}
{$81\text{ cm}$}
\loigiai{
Diện tích hình vuông: $9 \times 9 = 81\text{ (cm}^2)$.
\textbf{Đáp án của câu hỏi là B}
}
\end{ex}

\begin{ex} %[3H3V52]
Một hình chữ nhật có diện tích $72\text{ cm}^2$, chiều dài là $9\text{ cm}$. Chu vi của hình chữ nhật đó là
\choice
{$17\text{ cm}$}
{$36\text{ cm}$}
{\True $34\text{ cm}$}
{$32\text{ cm}$}
\loigiai{
Chiều rộng: $72 : 9 = 8\text{ cm}$. Chu vi: $(9 + 8) \times 2 = 34\text{ (cm)}$.
\textbf{Đáp án của câu hỏi là C}
}
\end{ex}

\begin{ex} %[3H3V52]
Một hình vuông có chu vi bằng $36\text{ cm}$. Diện tích của hình vuông đó là
\choice
{\True $81\text{ cm}^2$}
{$72\text{ cm}^2$}
{$64\text{ cm}^2$}
{$90\text{ cm}^2$}
\loigiai{
Cạnh hình vuông: $36 : 4 = 9\text{ cm}$. Diện tích: $9 \times 9 = 81\text{ (cm}^2)$.
\textbf{Đáp án của câu hỏi là A}
}
\end{ex}

\begin{ex} %[3H3V52]
Một mảnh đất hình chữ nhật có chiều rộng $8\text{ m}$, chiều dài gấp đôi chiều rộng. Diện tích mảnh đất đó là
\choice
{$48\text{ m}^2$}
{$64\text{ m}^2$}
{$256\text{ m}^2$}
{\True $128\text{ m}^2$}
\loigiai{
Chiều dài: $8 \times 2 = 16\text{ m}$. Diện tích: $16 \times 8 = 128\text{ (m}^2)$.
\textbf{Đáp án của câu hỏi là D}
}
\end{ex}

\begin{ex} %[3H3V52]
Nếu cạnh của một hình vuông tăng lên gấp 2 lần thì diện tích của hình vuông đó sẽ
\choice
{Tăng lên 2 lần}
{\True Tăng lên 4 lần}
{Tăng lên 8 lần}
{Không thay đổi}
\loigiai{
Diện tích hình vuông bằng cạnh $\times$ cạnh. Khi cạnh gấp đôi: $2 \times 2 = 4$ lần.
\textbf{Đáp án của câu hỏi là B}
}
\end{ex}

\begin{ex} %[3H3V52]
Một căn phòng hình chữ nhật có chiều dài $8\text{ m}$, chiều rộng $5\text{ m}$. Người ta lát nền bằng các viên gạch vuông cạnh $1\text{ m}$. Số viên gạch cần dùng là
\choice
{$26$ viên}
{$30$ viên}
{\True $40$ viên}
{$50$ viên}
\loigiai{
Diện tích nền: $8 \times 5 = 40\text{ m}^2$. Mỗi viên gạch $1\text{ m}^2$. Cần $40 : 1 = 40$ viên gạch.
\textbf{Đáp án của câu hỏi là C}
}
\end{ex}

\begin{ex} %[3H3V52]
Một tấm bìa hình vuông có diện tích $64\text{ cm}^2$. Chu vi của tấm bìa đó là
\choice
{\True $32\text{ cm}$}
{$16\text{ cm}$}
{$64\text{ cm}$}
{$24\text{ cm}$}
\loigiai{
Vì $8 \times 8 = 64$ nên cạnh hình vuông là $8\text{ cm}$. Chu vi: $8 \times 4 = 32\text{ (cm)}$.
\textbf{Đáp án của câu hỏi là A}
}
\end{ex}

\begin{ex} %[3H3V52]
Một tấm tôn hình chữ nhật có chiều dài $15\text{ cm}$, chiều rộng $9\text{ cm}$. Người ta cắt đi một góc hình vuông cạnh $4\text{ cm}$. Diện tích phần tôn còn lại là
\choice
{$135\text{ cm}^2$}
{$125\text{ cm}^2$}
{$120\text{ cm}^2$}
{\True $119\text{ cm}^2$}
\loigiai{
Diện tích ban đầu: $15 \times 9 = 135\text{ cm}^2$. Diện tích cắt đi: $4 \times 4 = 16\text{ cm}^2$. Còn lại: $135 - 16 = 119\text{ (cm}^2)$.
\textbf{Đáp án của câu hỏi là D}
}
\end{ex}

\begin{ex} %[3H3V52]
Người ta ghép 3 hình vuông bằng nhau cạnh $4\text{ cm}$ thành một hình chữ nhật. Diện tích hình chữ nhật đó là
\choice
{$36\text{ cm}^2$}
{\True $48\text{ cm}^2$}
{$64\text{ cm}^2$}
{$32\text{ cm}^2$}
\loigiai{
Diện tích mỗi hình vuông: $4 \times 4 = 16\text{ cm}^2$. Diện tích hình chữ nhật: $16 \times 3 = 48\text{ (cm}^2)$.
\textbf{Đáp án của câu hỏi là B}
}
\end{ex}

\begin{ex} %[3H3V52]
Một hình chữ nhật có chu vi $28\text{ cm}$, chiều dài là $8\text{ cm}$. Diện tích của hình chữ nhật đó là
\choice
{$40\text{ cm}^2$}
{$36\text{ cm}^2$}
{\True $48\text{ cm}^2$}
{$56\text{ cm}^2$}
\loigiai{
Nửa chu vi: $28 : 2 = 14\text{ cm}$. Chiều rộng: $14 - 8 = 6\text{ cm}$. Diện tích: $8 \times 6 = 48\text{ (cm}^2)$.
\textbf{Đáp án của câu hỏi là C}
}
\end{ex}

\begin{ex} %[3H3V52]
So sánh diện tích hình chữ nhật kích thước $6\text{ cm} \times 6\text{ cm}$ (hình vuông) và hình chữ nhật kích thước $8\text{ cm} \times 4\text{ cm}$. Khẳng định đúng là
\choice
{\True Hình vuông có diện tích lớn hơn $4\text{ cm}^2$}
{Hình chữ nhật lớn hơn}
{Hai hình có diện tích bằng nhau}
{Không so sánh được}
\loigiai{
Diện tích hình vuông: $6 \times 6 = 36\text{ cm}^2$. Diện tích hình chữ nhật: $8 \times 4 = 32\text{ cm}^2$. Lớn hơn: $36 - 32 = 4\text{ cm}^2$.
\textbf{Đáp án của câu hỏi là A}
}
\end{ex}

\textbf{Phần 2. Trắc nghiệm đúng sai}
\setcounterpageref{ex}{0}

\begin{ex} %[3H3V52]
Xét tính đúng sai của công thức tính diện tích hình phẳng:
\choiceTF
{\True Diện tích hình chữ nhật bằng chiều dài nhân với chiều rộng (cùng đơn vị đo).}
{\True Diện tích hình vuông bằng độ dài một cạnh nhân với chính nó.}
{Một hình vuông cạnh $5\text{ cm}$ có diện tích là $20\text{ cm}^2$.}
{Đơn vị đo chu vi và diện tích là như nhau.}
\loigiai{
\begin{itemize}
    \item Ý a) Đúng: Công thức chuẩn.
    \item Ý b) Đúng: Công thức chuẩn.
    \item Ý c) Sai: Sai vì $5 \times 5 = 25\text{ cm}^2$ ($20\text{ cm}$ là chu vi).
    \item Ý d) Sai: Sai vì chu vi đo bằng $\text{cm}$, diện tích đo bằng $\text{cm}^2$.
\end{itemize}
\begin{itemize}
    \item a) \textbf{Đúng}.
    \item b) \textbf{Đúng}.
    \item c) \textbf{Sai}.
    \item d) \textbf{Sai}.
\end{itemize}
}
\end{ex}

\begin{ex} %[3H3V52]
Một hình chữ nhật có chiều dài $10\text{ cm}$, chiều rộng $5\text{ cm}$:
\choiceTF
{\True Chu vi hình chữ nhật bằng $30\text{ cm}$.}
{\True Diện tích hình chữ nhật bằng $50\text{ cm}^2$.}
{Nếu tăng chiều rộng lên $5\text{ cm}$ để thành hình vuông thì diện tích không đổi.}
{\True Hình chữ nhật có thể chia thành 2 hình vuông cạnh $5\text{ cm}$.}
\loigiai{
\begin{itemize}
    \item Ý a) Đúng: $(10 + 5) \times 2 = 30\text{ cm}$.
    \item Ý b) Đúng: $10 \times 5 = 50\text{ cm}^2$.
    \item Ý c) Sai: Sai vì hình vuông mới có diện tích $10 \times 10 = 100\text{ cm}^2$.
    \item Ý d) Đúng: $10 : 5 = 2$ hình vuông.
\end{itemize}
\begin{itemize}
    \item a) \textbf{Đúng}.
    \item b) \textbf{Đúng}.
    \item c) \textbf{Sai}.
    \item d) \textbf{Đúng}.
\end{itemize}
}
\end{ex}

\textbf{Phần 3. Trả lời ngắn}
\setcounterpageref{ex}{0}

\begin{ex} %[3H3V52]
Tính diện tích hình chữ nhật có chiều dài $14\text{ cm}$, chiều rộng $6\text{ cm}$.
\shortans{84}
\loigiai{
\textbf{Đáp số:} $84$\\
$14 \times 6 = 84\text{ (cm}^2)$.
}
\end{ex}

\begin{ex} %[3H3V52]
Một hình vuông có cạnh $8\text{ cm}$. Tính diện tích hình vuông đó theo xăng-ti-mét vuông.
\shortans{64}
\loigiai{
\textbf{Đáp số:} $64$\\
$8 \times 8 = 64\text{ (cm}^2)$.
}
\end{ex}

\begin{ex} %[3H3V52]
Một hình chữ nhật có diện tích $45\text{ cm}^2$ và chiều rộng $5\text{ cm}$. Tính chu vi hình chữ nhật đó theo xăng-ti-mét.
\shortans{28}
\loigiai{
\textbf{Đáp số:} $28$\\
Chiều dài: $45 : 5 = 9\text{ cm}$. Chu vi: $(9 + 5) \times 2 = 28\text{ (cm)}$.
}
\end{ex}

\begin{ex} %[3H3V52]
Hình vuông có chu vi $28\text{ cm}$. Tính diện tích của hình vuông đó theo xăng-ti-mét vuông.
\shortans{49}
\loigiai{
\textbf{Đáp số:} $49$\\
Cạnh: $28 : 4 = 7\text{ cm}$. Diện tích: $7 \times 7 = 49\text{ (cm}^2)$.
}
\end{ex}

\begin{ex} %[3H3V52]
Một hình chữ nhật có chiều dài $15\text{ cm}$, chiều rộng bằng $\dfrac{1}{3}$ chiều dài. Tính diện tích hình chữ nhật đó.
\shortans{75}
\loigiai{
\textbf{Đáp số:} $75$\\
Chiều rộng: $15 : 3 = 5\text{ cm}$. Diện tích: $15 \times 5 = 75\text{ (cm}^2)$.
}
\end{ex}

\begin{ex} %[3H3V52]
Cần bao nhiêu viên gạch men hình vuông cạnh $2\text{ dm}$ để lát kín một căn phòng diện tích $20\text{ m}^2$ (biết $1\text{ m}^2 = 25$ viên gạch)?
\shortans{500}
\loigiai{
\textbf{Đáp số:} $500$\\
Số viên gạch: $20 \times 25 = 500$ viên.
}
\end{ex}
